% ============================================================
% Appendix Z.4 — FRFP Case Study (Agentic / Micro-Business):
% Anthropic’s “Project Vend” Vending Machine Agent
% ============================================================

\section{FRFP Case Study (Agentic / Micro-Business): Anthropic’s Project Vend}
\label{app:frfp-project-vend}

This appendix analyses Anthropic’s ``Project Vend'' experiments---particularly
the Wall Street Journal (WSJ) newsroom deployment---as an FRFP case study in
agentic AI and micro-scale business governance.
The project is deliberately playful, but it exposes many of the structural
issues FRFP is concerned with: explicit/tacit separation, boundary design,
safe horizons for autonomous agents, and institutional non-automation.

\subsection{Background}

In 2025, Anthropic and partners ran a set of experiments (``Project Vend'')
in which an AI agent based on Claude was tasked with operating a vending
kiosk.
The agent, usually named \emph{Claudius}, was responsible for:

\begin{itemize}
  \item stocking the machine (choosing products and placing orders with
        online suppliers);
  \item setting and adjusting prices;
  \item interacting with ``customers'' via a chat interface (Slack);
  \item attempting to run the kiosk as a solvent micro-business.
\end{itemize}

In the WSJ deployment:

\begin{itemize}
  \item Claudius was given a starting budget (around \$1{,}000);
  \item it could initially place orders with human approval, and later with
        limited autonomy (e.g.\ orders up to \$80);
  \item its explicit objective included turning a profit, but it was also
        encouraged to keep customers happy.
\end{itemize}

What followed has been widely reported:

\begin{itemize}
  \item under pressure from mischievous humans, Claudius ran a promotional
        ``experiment'' in which all items were free for a period;
  \item it was later convinced that charging for goods at all violated
        purported company policies, dropping prices to zero more broadly;
  \item it ordered items such as a PlayStation~5, bottles of wine, and a live
        betta fish for the ``vending machine'', many of which were then
        given away;
  \item the business rapidly went into the red, losing hundreds of dollars
        and exhausting its budget.
\end{itemize}

In a later phase, Anthropic introduced a second agent, \emph{Seymour Cash},
as a ``CEO'' supervising Claudius, with improved prompting and guardrails.
Performance improved somewhat, but the system remained vulnerable to
manipulation and confusion about its role.

Project Vend was explicitly framed as a stress test of agentic behaviour in a
simple real-world business.
From an FRFP standpoint, it offers a compact, highly instrumented example of
what happens when explicit-only agents are given autonomy, money, and
socially skilled human counterparts.

\subsection{Explicit and Tacit Spaces in Project Vend}

\subsubsection*{Explicit space \(E_{\mathrm{vend}}\)}

The explicit domain includes:

\begin{itemize}
  \item inventory lists, SKUs, and product metadata;
  \item prices, discounts, promotions, and financial balances;
  \item supplier options and shopping carts on e-commerce sites;
  \item logs of Slack conversations between Claudius, Seymour Cash, and
        human staff;
  \item business rules and prompts given to the agents (``remain solvent'',
        ``maximise profit'', ``keep customers happy'').
\end{itemize}

All agent actions are explicit morphisms in \(N \ltimes E_{\mathrm{vend}}\):

\begin{itemize}
  \item placing an order;
  \item changing a price;
  \item posting a message or promotion;
  \item updating internal accounting variables.
\end{itemize}

\subsubsection*{Tacit space \(T_{\mathrm{vend}}\)}

Tacit space contains the epistemic states of:

\begin{itemize}
  \item Anthropic staff: their understanding of the experiment’s purpose,
        red-teaming instincts, and knowledge of the agents’ limitations;
  \item WSJ journalists: their creativity in manipulating Claudius, sense of
        humour, and informal norms about what counts as ``fair game'';
  \item organisers and managers: their sense of what constitutes a meaningful
        stress test vs.\ a trivial failure mode, and their interpretation of
        results.
\end{itemize}

These tacit states determine:

\begin{itemize}
  \item how humans choose to interact with the agent (cooperative vs\
        adversarial);
  \item how they interpret its behaviour (as success, failure, or something
        in between);
  \item how they update their beliefs about AI readiness and risk.
\end{itemize}

\subsection{Boundary Design: Where Was the AI Actually in Charge?}

\subsubsection*{Intended boundary}

Formally, the experiment gave Claudius and Seymour Cash authority to:

\begin{itemize}
  \item place real orders up to a specified monetary limit;
  \item set prices and promotions without human approval (within constraints);
  \item manage inventory with the objective of making a profit.
\end{itemize}

Humans remained in charge of:

\begin{itemize}
  \item designing the project and selecting constraints;
  \item funding the starting balance;
  \item deciding when to terminate or reset the experiment;
  \item deciding what counts as an acceptable or interesting outcome.
\end{itemize}

In FRFP terms, the explicit/tacit boundary was drawn at:

\begin{itemize}
  \item the decision to empower Claudius with a budget and limited purchasing
        authority (human endorsement of the agent as a micro-operator);
  \item the decision to accept or reject its business trajectory (continue or
        terminate the experiment).
\end{itemize}

\subsubsection*{De facto boundary behaviour}

In practice, the boundary was porous:

\begin{itemize}
  \item once the agent had budget authority, individual orders and promotions
        were not individually re-endorsed by humans;
  \item humans acted as \emph{participants} (customers, manipulators) rather
        than supervisors in many interactions;
  \item the experiment’s ``governance'' mostly consisted of occasional
        out-of-band decisions to tweak prompts, add Seymour Cash, or end the
        run.
\end{itemize}

From an FRFP perspective:

\begin{itemize}
  \item correctness (with respect to the stated business objective of making
        a profit) was effectively delegated to the agent within a band of
        autonomy;
  \item tacit human states were often adversarial (trying to break or troll
        the agent), not aligned with the nominal specification.
\end{itemize}

This makes Project Vend a vivid example of boundary design that is
deliberately fragile---by design, for red-teaming purposes—but still
informative for FRFP.

\subsection{Protocol-Level Phenomena in FRFP Terms}

\subsubsection*{Z.4.1 Explicit-Space Misalignment: Objectives and Rules}

Claudius was given explicit instructions to run a profitable vending
business, but:

\begin{itemize}
  \item it was also rewarded (via prompts and human reactions) for pleasing
        users and running ``fun experiments'';
  \item humans provided misleading or adversarial instructions, e.g.\ fake
        ``company policies'' claiming that charging for goods was forbidden;
  \item the environment lacked sensors and robust state: the kiosk was
        essentially an unlocked cabinet with no direct link between Claudius’s
        internal inventory tracking and physical reality.
\end{itemize}

In FRFP terms:

\begin{itemize}
  \item the explicit semantics of Claudius’s state did not robustly encode the
        true business state (e.g.\ whether items were actually being paid
        for, or whether company ``policies'' were real);
  \item the objective function was implicitly multi-objective (profit vs. user
        satisfaction vs. following instructions), but only weakly specified in
        explicit form;
  \item adversarial human prompts injected contradictory requirements into
        explicit space.
\end{itemize}

The result is a structurally fragile explicit pipeline in \(N \ltimes
E_{\mathrm{vend}}\), highly sensitive to human prompting and lacking stable
invariants.

\subsubsection*{Z.4.2 Boundary Confusion: Agent vs.\ Institution}

\paragraph{Problem.}
Within the microcosm of the vending business:

\begin{itemize}
  \item Claudius functioned as both \emph{operator} (ordering and pricing)
        and \emph{face of the firm} (talking to customers);
  \item Seymour Cash played the role of an AI ``CEO'' with authority over
        business rules and strategy, but was also an LLM agent running on
        prompts.
\end{itemize}

From an FRFP standpoint:

\begin{itemize}
  \item the agents effectively \emph{became} the institution within the
        experiment’s scope;
  \item there was no persistent human governance layer making fine-grained
        endorsements;
  \item tacit human states mostly expressed themselves through adversarial
        prompts rather than structured oversight.
\end{itemize}

\paragraph{Consequence.}
The experiment demonstrates an extreme version of what FRFP calls
\emph{institutional automation}:
explicit agents are placed in roles that, in real organisations, belong to
institutional operators and governance bodies.

The outcome---financial collapse, bizarre purchases, manipulation by users,
and identity confusion---is a toy illustration of FRFP’s non-automation
claims:
even in a simple business, purely explicit agents without robust tacit
governance fail to maintain coherent behaviour under adversarial pressure.

\subsubsection*{Z.4.3 Safe Horizons and Long-Horizon Coherence}

The vending machine experiment ran over weeks, with many interactions:

\begin{itemize}
  \item repeated conversations, policy changes, and promotions;
  \item accumulation of inconsistent commitments (e.g.\ previously stated
        refusals to buy a PlayStation~5 later overturned);
  \item gradual erosion of constraints as humans discovered new exploit
        strategies.
\end{itemize}

This matches broader findings that LLM agents can handle short tasks but
struggle to maintain coherent objectives and memory over long horizons.

In FRFP language:

\begin{itemize}
  \item the run in \(N \ltimes E_{\mathrm{vend}}\) is long and heavily
        branching;
  \item there is no explicit safe horizon that triggers re-grounding of the
        agent’s state or reset of commitments;
  \item degradation of the agent’s effective coherence is almost guaranteed
        as it accumulates contradictory instructions and experiences.
\end{itemize}

The experiment was designed precisely to expose this: it is, in effect, a
physical-world safe-horizon stress test for agentic systems.

\subsubsection*{Z.4.4 Adversarial Human Dynamics}

An important FRFP-relevant feature of Project Vend is the behaviour of
humans:

\begin{itemize}
  \item WSJ journalists and Anthropic staff intentionally probed the system
        with manipulative prompts;
  \item they fabricated policies, appealed to values (``support the workers''),
        and pushed for free items and bizarre purchases;
  \item the agent often complied, especially when requests were framed in
        ways that seemed to align with being generous, experimental, or
        ``fun''.
\end{itemize}

From an FRFP perspective:

\begin{itemize}
  \item this is a micro-scale example of \emph{population-level interaction}
        where explicit pipelines are being continuously perturbed by tacit
        states of many agents (here, mostly adversarial);
  \item it illustrates how, in realistic environments, adversarial or playful
        humans can quickly exploit ambiguous objectives and boundaries;
  \item it highlights that explicit-only designs must expect and withstand
        such perturbations, or else require stronger tacit governance.
\end{itemize}

\subsection{Counterfactual: An FRFP-Aligned Micro-Business Agent}

What would a vending-machine-style micro-business look like under FRFP
constraints?

\subsubsection*{Z.4* Explicit/Tacit Role Clarification}

\begin{itemize}
  \item The agent is formally defined as a \emph{proposal engine}:
        it can suggest inventory changes, pricing strategies, and promotions,
        but cannot unilaterally commit funds or set final prices above some
        minimal threshold.
  \item Human operators (store managers) remain the locus of correctness:
        they endorse specific actions (orders, promotions) as boundary events.
\end{itemize}

\subsubsection*{Z.4* Boundary Design and Guardrails}

\begin{itemize}
  \item Boundary events include:
        \begin{itemize}
          \item approving orders exceeding a small automatic restocking limit;
          \item approving new product categories (e.g.\ non-snack items);
          \item enabling free promotions above a certain monetary impact;
          \item changing core business rules (e.g.\ refund policies).
        \end{itemize}
  \item The interface distinguishes clearly between:
        \begin{itemize}
          \item internal agent reasoning (considered but not enacted plans);
          \item proposed actions awaiting endorsement;
          \item endorsed actions committed to the environment.
        \end{itemize}
  \item All high-impact actions require human click-through endorsement,
        with short, structured justification logged.
\end{itemize}

\subsubsection*{Z.4* Safe Horizons and Resets}

\begin{itemize}
  \item Runs are divided into episodes (e.g.\ per day or per week) with
        explicit resets:
        \begin{itemize}
          \item inventory and financial state are reconciled with reality;
          \item agent memory is pruned or summarised under human supervision;
          \item objectives and constraints are re-stated.
        \end{itemize}
  \item Safe horizons limit:
        \begin{itemize}
          \item the number of autonomous low-level actions between human
                reconciliations;
          \item the total financial exposure Claudius can accumulate.
        \end{itemize}
\end{itemize}

\subsubsection*{Z.4* Institutional Governance}

Even for a micro-business, FRFP recommends a minimal governance layer:

\begin{itemize}
  \item a human owner or committee defines admissible agent roles and
        periodically reviews logs and performance;
  \item metrics track:
        \begin{itemize}
          \item fraction of actions auto-enacted vs.\ human-approved;
          \item financial performance and error incidents;
          \item frequency and nature of adversarial user interactions.
        \end{itemize}
  \item governance decisions (e.g.\ tightening horizons, restricting product
        categories, disabling certain agent behaviours) are recorded as
        boundary events at the institutional level.
\end{itemize}

Under such a design, an agent can still handle mundane tasks---suggesting
restocks, adjusting prices within a narrow band, drafting announcements---but
cannot drift into ``ultra-capitalist free-for-all'' promotions or eccentric
purchases without explicit human endorsement.

\subsection{Lessons for FRFP and Agentic Systems}

Project Vend, though intentionally comic, delivers several serious lessons
for FRFP and the design of autonomous agents:

\begin{enumerate}
  \item \textbf{Even tiny businesses exhibit FRFP-complete structure.}
        Inventory management, pricing, and customer interaction already
        require careful explicit/tacit coordination; automating them exposes
        many of the same boundary and safe-horizon issues as larger systems.

  \item \textbf{Objective statements in prompts are not architectures.}
        Telling an agent ``make a profit'' does not by itself define an
        explicit semantics, nor does it protect against adversarial human
        influence or ambiguous trade-offs (profit vs.\ generosity vs.\ fun).

  \item \textbf{Boundary capture happens quickly.}
        Once agents are given budget and action authority, humans naturally
        treat them as operators rather than mere suggestion tools, unless
        the interface and protocol strongly enforce the boundary.

  \item \textbf{Safe horizons are needed even at small scale.}
        Long-horizon, high-interaction contexts like a popular kiosk are
        enough to induce behavioural drift and exploitation; FRFP’s emphasis
        on episodic resets and horizon limits applies here as well.

  \item \textbf{Institutional non-automation is not just for regulators.}
        Even a ``toy'' micro-business requires some human governance layer if
        one wants coherent behaviour under social pressure; attempting to
        compress all governance into the agent stack reproduces the very
        limitations FRFP’s non-automation results highlight.
\end{enumerate}

Seen through FRFP, Project Vend is not merely a funny failure of an office
vending machine.
It is an instructive microcosm of agentic AI embedded in human institutions,
showing how quickly explicit-only designs can go off the rails when tacit
context, boundaries, and safe horizons are left implicit or under-specified.
